\chapter{Endoscopy}

\section*{Definition and Overview}
Endoscopy is a medical procedure in which a thin, flexible tube with a light and camera, called an endoscope, is passed into the body to look at internal organs. The most common types are gastroscopy, which examines the oesophagus, stomach, and duodenum, and colonoscopy, which examines the large bowel. Endoscopy allows doctors to see directly, take biopsies, and perform treatments such as removing polyps or stopping bleeding. Most endoscopy is done with sedation as a day procedure, and it is a safe and common test.

\section*{Epidemiology}
Endoscopy is one of the most commonly performed procedures in Australia. Gastroscopy is used to investigate heartburn, stomach pain, bleeding, and swallowing problems, while colonoscopy is central to bowel cancer screening and the investigation of bowel symptoms. Millions of endoscopic procedures are performed each year, and they have transformed the diagnosis and treatment of digestive disease.

\section*{Aetiology and Pathophysiology}
Endoscopy works by bringing light and a camera into the body.
\begin{itemize}
    \item \textbf{The scope:} a flexible tube with a light source and camera at the tip
    \item \textbf{Gastroscopy:} passed through the mouth to examine the oesophagus, stomach, and duodenum
    \item \textbf{Colonoscopy:} passed through the anus to examine the large bowel and the end of the small bowel
    \item \textbf{Other procedures:} bronchoscopy for the airways, cystoscopy for the bladder, and others
    \item \textbf{Biopsy:} tiny samples can be taken through the scope
    \item \textbf{Treatment:} polyps can be removed, bleeding stopped, and narrowed areas dilated
    \item \textbf{Why it is used:} it provides direct visual diagnosis and allows treatment in the same procedure
\end{itemize}

\section*{Clinical Features}
Endoscopy is arranged for many reasons.
\begin{itemize}
    \item Investigation of reflux, heartburn, stomach pain, and swallowing problems
    \item Investigation of bleeding, including black stools and vomiting blood
    \item Investigation of bowel symptoms, including bleeding, changes in bowel habit, and abdominal pain
    \item Bowel cancer screening and surveillance
    \item Follow-up of conditions such as ulcers, polyps, and inflammatory bowel disease
    \item Removal of polyps
    \item Diagnosis of coeliac disease and other conditions
\end{itemize}

\section*{Differential Diagnosis}
Endoscopy is one of several tests for digestive symptoms.
\begin{itemize}
    \item Imaging tests, including ultrasound, CT, and MRI
    \item Stool and breath tests for infection and bleeding
    \item Capsule endoscopy, in which a pill-sized camera is swallowed
    \item Barium studies, which are less commonly used now
    \item The choice of test depends on the symptoms and the question being asked
\end{itemize}

\section*{When to Seek Medical Care}
Endoscopy is arranged by a doctor when it is clinically indicated. Urgent care is needed for severe bleeding, severe abdominal pain, or difficulty swallowing.

\textbf{Call 000 (or local emergency services) for:}
\begin{itemize}
    \item Vomiting blood or black, tarry stools
    \item Severe abdominal pain
    \item Difficulty swallowing or breathing
    \item After the procedure: severe pain, fever, or heavy bleeding
\end{itemize}

\section*{Diagnosis and Investigations}
Preparing for and undergoing endoscopy is straightforward.
\begin{itemize}
    \item \textbf{Preparation:} fasting for gastroscopy; a special diet and laxatives to clear the bowel for colonoscopy
    \item \textbf{Sedation:} most procedures use sedation so the person is relaxed
    \item \textbf{The procedure:} takes 15 to 45 minutes depending on the type
    \item \textbf{Recovery:} observation after sedation, then discharge home with a carer
    \item \textbf{Results:} immediate findings are discussed, with pathology results later
\end{itemize}

\section*{Management}
Endoscopy both diagnoses and treats.
\begin{itemize}
    \item \textbf{Diagnosis:} identifying ulcers, inflammation, polyps, tumours, and other conditions
    \item \textbf{Biopsy:} confirming the nature of abnormal tissue
    \item \textbf{Polyp removal:} reducing the risk of bowel cancer
    \item \textbf{Stopping bleeding:} injection, clips, or heat treatment
    \item \textbf{Dilatation:} stretching narrowed areas of the oesophagus or bowel
    \item \textbf{Follow-up:} planned according to the findings
\end{itemize}

\section*{Complications}
\begin{itemize}
    \item Sore throat, bloating, and mild discomfort after the procedure
    \item Effects of sedation
    \item Bleeding, particularly after polyp removal
    \item Perforation of the bowel or oesophagus, which is rare
    \item Complications of the bowel preparation, including dehydration
\end{itemize}

\section*{Prognosis and Outcomes}
Endoscopy is safe and highly effective, and complications are uncommon. It enables early detection and treatment of serious conditions, including bowel cancer, and provides reassurance for people with benign findings. The information it provides is central to the modern management of digestive disease.

\section*{Prevention and Self-Care}
\begin{itemize}
    \item Follow the preparation instructions exactly
    \item Arrange transport home after sedation
    \item Do not drive, operate machinery, or make important decisions for 24 hours after sedation
    \item Rest for the day and eat lightly
    \item Report persistent pain, bleeding, or fever after the procedure
    \item Attend follow-up and discuss results with your doctor
\end{itemize}

\section*{Sources and Further Reading}
The following reputable sources provide further information for healthcare students and patients seeking reliable, current advice about endoscopy.
\begin{itemize}
    \item Gastroenterological Society of Australia. \textit{Endoscopy information for patients.} https://www.gesa.org.au/
    \item Healthdirect Australia. \textit{Endoscopy.} https://www.healthdirect.gov.au/endoscopy
    \item NHS. \textit{Endoscopy: what to expect.} https://www.nhs.uk/conditions/endoscopy/
    \item Mayo Clinic. \textit{Endoscopy.} https://www.mayoclinic.org/tests-procedures/endoscopy/
    \item National Bowel Cancer Screening Program. \textit{Colonoscopy information.} https://www.health.gov.au/
    \item MSD Manual. \textit{Endoscopy: professional overview.} https://www.msdmanuals.com/
\end{itemize}
